\documentclass{beamer}
\usepackage{etex}
\usepackage{lmodern}
\usepackage[english]{babel}
\usepackage[T1]{fontenc}
\usepackage[utf8]{inputenc}
\usepackage{pst-sigsys} 
\usepackage{amsmath,amsfonts,amssymb}
\usepackage{pstricks-add} 
\usepackage{ragged2e}
\usepackage{graphicx}
\usepackage{ulem}
\usepackage{fontawesome}
\setbeamertemplate{navigation symbols}{}  
 
\usetheme{Darmstadt}
\setbeamertemplate{footline}{\insertframenumber/\inserttotalframenumber}
\title{M1 IEAP - BTI/FH/IEMH\\FIEA11CM : Analyse et Traitement du Signal}
\author{Flavy ROSEREN\\Martin EGIZIANO\\Frank BULOUP}
\institute{Aix Marseille Université\\Institut des Sciences du Mouvement}
\date{}

\setbeamertemplate{footline} 
{  
	\begin{beamercolorbox}[ht=2.5ex,dp=1.125ex,%
      leftskip=.3cm,rightskip=.3cm plus1fil]{title in head/foot}%
      {\usebeamerfont{title in head/foot}\insertshorttitle} \hfill    
      \insertframenumber / \inserttotalframenumber%
    \end{beamercolorbox}%
%     \begin{beamercolorbox}[colsep=1.5pt]{lower separation line foot}
%     \end{beamercolorbox} 
}

\newcounter{exampleBlockCounter}
\setcounter{exampleBlockCounter}{1}

\begin{document} 
 
\begin{frame}[plain] 
	\titlepage 
	\center{\includegraphics[scale=0.75]{images/by-nc-sa.eps}}
	\vspace{1cm}
	\includegraphics[scale=0.6]{images/LogoAMU.png}\hspace*{2cm}
	\includegraphics[scale=0.2]{images/LogoCNRS.eps}\hspace*{2cm}
	\includegraphics[scale=0.1]{images/LogoISM.eps}
\end{frame}

\begin{frame}{Première Partie}
	\tableofcontents
\end{frame}

\section{Signal Discret - Aspect temporel}
\subsection{Discrétisation : lien avec l'échantillonnage}
\begin{frame}
\begin{block}{Discrétisation}
\justifying
Opération consistant à remplacer des relations portant sur des fonctions
continues par un nombre fini de relations algébriques 
portant sur les valeurs prises par ces fonctions en un nombre fini de points
 de leur ensemble de définition (Def. Larousse).
 \smallbreak
 \center
 \textcolor{red}{Lien évident avec l'échantillonnage, l'acquisition de données}
\end{block}
\pause
\begin{exampleblock}{Exercice \Roman{exampleBlockCounter} - Exemple}
\justifying
Soit la fonction $f(t) = Acos(\omega t)$. Prendre les valeurs de $f(t)$ tous les
$nT_{e}$ revient à discrétiser f(t). On parle alors de fonction discrète ou
plutôt de signal discret et on note $f(n) = Acos(\omega nT_{e})$
\end{exampleblock}
\end{frame}
\stepcounter{exampleBlockCounter}
\begin{frame}
	\begin{exampleblock}{Exercice \Roman{exampleBlockCounter} - Application sous Python}
		\begin{enumerate}\justifying
			  \item Etudier la fonction \textit{wavread} dans l'aide Python
			  \item Charger le signal contenu dans le fichier \textit{mi2\_bute\_ros.wav} ainsi que la fréquence d'échantillonnage d'enregistrement $F_{s}$
			  \item Ecouter le fichier (fonction \textit{soundsc}).  Que se passe-t-il si on multiplie ou divise $F_{s}$ par 2 ?
			  \item Créer un vecteur temporel correspondant au signal. Tracer le signal en fonction du temps.
			  \item En utilisant les outils de zoom et le curseur, estimer la fréquence fondamentale de la note jouée.
			  \item Sur un second graphe (fonction \textit{subplot}) tracer le signal pour une fréquence d'échantillonnage deux fois plus élevée. Vérifier l'observation faite à la question 3.   
		\end{enumerate}
	\end{exampleblock}
\end{frame}
\stepcounter{exampleBlockCounter}

\subsection[Signaux élémentaires]{Quelques signaux discrets élémentaires}
\begin{frame}
\begin{block}{L'impulsion unitaire}
\begin{columns}[t]
	\begin{column}{5cm}
	\smallbreak 
		\[
			\delta (n) = 
			\begin{cases}
			1 & \text{si } n = 0 \\
			0 & \text{ailleurs}
			\end{cases}
		\]
	\end{column}
	\begin{column}{5cm}
		\begin{figure}
			\begin{pspicture}[showgrid=false](-2,-.5)(3,1.5)
				\psaxeslabels{->}(0,0)(-2,0)(3,1.5){$n$}{$\delta (n)$}
				\psstem[stemhead=o, stemtag](-2,1){0, 0, 1 ,0 , 0}  
			\end{pspicture}
		\end{figure}
	\end{column}
\end{columns}
\end{block}
\pause
\begin{block}{L'échelon unitaire}
\begin{columns}[t]
	\begin{column}{5cm}
	\smallbreak 
		\[
		u(n) =
		\begin{cases}
		1 & \text{si } n \geq 0 \\
		0 & \text{ailleurs}
		\end{cases}  
		\]
	\end{column}
	\begin{column}{5cm}
		\begin{figure}
			\begin{pspicture}[showgrid=false](-2,-.5)(3,1.5)
				\psaxeslabels{->}(0,0)(-2,0)(3,1.5){$n$}{$u(n)$}
				\psstem[stemhead=o, stemtag](-2,1){0, 0, 1 , 1, 1}  
			\end{pspicture}
		\end{figure}
	\end{column}
\end{columns}
\end{block}
\end{frame}

\begin{frame}
\begin{block}{La porte}
\begin{columns}[t]
	\begin{column}{5cm}
	\smallbreak 
		\[
		\Pi_{N}(n)=
		\begin{cases}
		1 & \text{si } 0 \leq n \leq N-1 \\
		0 & \text{ailleurs}
		\end{cases}
		\]
	\end{column}
	\begin{column}{5cm}
		\begin{figure}
			\begin{pspicture}[showgrid=false](-2,-.5)(3,1.5)
				\psaxeslabels{->}(0,0)(-2,0)(3,1.5){$n$}{$\Pi_{2}(n)$}
				\psstem[stemhead=o, stemtag](-2,1){0, 0, 1 , 1, 0}  
			\end{pspicture}
		\end{figure}
	\end{column}
\end{columns}
\end{block}
\pause
\begin{alertblock}{Remarques}
\[
u(n) = \sum_{k=0}^{+\infty} \delta (n-k)
\] 
\[
\Pi_{N}(n) = \sum_{k=0}^{+\infty} \delta (n-k) - \sum_{l=N}^{+\infty}
\delta(n-l)
\] 
\end{alertblock} 
\end{frame}

\begin{frame}
\begin{block}{Les fonctions cosinus et sinus discrètes}
	\[
		s_1(n) = cos(2\pi f \,  nT_e)\text{ et }s_2(n) = sin(2\pi f \, nT_e)
	\]
\end{block}

\begin{alertblock}{Remarque}
\justifying
Les variables $n$, $k$ ou $l$ utilisées précédement sont implicitement
équivalentes à $nT_e$, $kT_e$ ou $lT_e$, $T_e$ étant la période d'échantillonnage 
\end{alertblock} 
\pause
\begin{exampleblock}{Exercice \Roman{exampleBlockCounter} - Signaux élémentaires}\justifying 
\itemize
	\item Tracer les signaux élémentaires sous Python %\sout{(voir éventuellement les
	%fonctions \textbf{square}, \textbf{sawtooth} et \textbf{chirp})}
	%\item Lire les fichiers de données MOUT104
\end{exampleblock}
\end{frame}  
\stepcounter{exampleBlockCounter}
\begin{frame} 
	\begin{exampleblock}{Exercice \Roman{exampleBlockCounter} - Transition vers l'aspect fréquentiel}
	\justifying Cet exercice propose de mettre en oeuvre la synthèse d'un signal
	par composition de sinusoïde
		\begin{enumerate}
		 	\item Créer un script nommé \textit{Exercice6.m} contenant "clear all;
		 	close all; clc;" comme premières commandes
			\item Créer un vecteur temporel dans l'intervalle $[0;2]$ avec un pas de
			  $0.001s$ 
			\item Tracer le signal discret $s(n)$ construit sur la composition suivante :
			\[ 
			s(t) = \frac{4}{\pi}[sin(2\pi t) + \frac{1}{3}sin(2\pi 3 t) +
			\frac{1}{5}sin(2\pi 5 t)]
			\] 
			\item Modifier le script pour sommer cinquante termes de la forme :
			\[
			\frac{4}{(2k+1)\pi}sin[2\pi (2k+1)t]\text{ $k$ allant de $0$ à  $49$}
			\] 
		\end{enumerate}
	\end{exampleblock}
\end{frame}
\stepcounter{exampleBlockCounter}
\section{Signal discret - Aspect fréquentiel}
\subsection{Série de Fourier} 
\begin{frame}
\begin{block}{Généralisation}
\justifying L'exemple de l'exercice précédent peut se généraliser  :
\textcolor{red}{presque tout \textbf{signal continu} périodique peut être
synthétisé par sommation de sinus et/ou de cosinus}. C'est ce que l'on appelle la \textcolor{red}{Série de Fourier}. Un signal $s(t)$ de
période $T_0 = \frac{1}{f_0}$ peut alors s'écrire de la façon suivante :
\begin{align*}
s(t) = S(0) + \sum_{k=1}^{+\infty} 2\rho (k) cos[2\pi kf_0t+\phi (k)]
\end{align*}
\end{block}
\pause
\begin{alertblock}{Remarques}
	\itemize
	\item $\rho(k)$ et $\phi(k)$ sont respectivement des amplitudes et des phases
	\pause
	\item $S(0)$ est la valeur moyenne, c'est toujours un réel pur
	\pause
	\item $f_0$ est la fréquence fondamentale 
	\pause
	\item Les fréquences des cosinus sont des multiples entiers de $f_0$. Ce sont
	les harmoniques.
\end{alertblock}
\end{frame}

\end{document}
